\section{Limitations}
Our evidence is specific to an ARK-style two-tool interface, \stark graphs, and open-weight \qwen controllers served locally. We evaluate query-derived profiles rather than real user memories, so the experiments should not be read as a complete personalized-retrieval benchmark. The held-out evaluation is a subset rather than the full test split, and the constrained-budget probe is diagnostic rather than a tuned production system. These constraints bound our claim that \tb shows that a plausible profile/trace-conditioned controller can degrade adaptive graph retrieval and that trace diagnostics explain the degradation. It, however, does not prove that all adaptive or personalized retrieval interventions will fail.

\section{Future work}
The next step is to replace prompt-level intervention with policy learning. The traces collected here define the state variables: candidate coverage, repeated calls, zero-result calls, relation coverage, budget use, and rank movement. A learned controller can use these variables to choose among global search, neighborhood exploration, relation relaxation, type tightening, reranking, and stopping. Personalization should enter this policy selectively. For example, a user profile may be useful for query anchoring or final reranking but harmful during relation expansion if it distracts from graph constraints. A validation-trained router that decides when to activate profile conditioning is a natural bridge from this diagnostic study to robust personalized retrieval.

\paragraph{A concrete policy-learning formulation.}
A direct extension is to train a small controller or router over the trace state rather than rely on unconditional textual hints. The training instance would consist of the query profile, current candidate summary, trace counters, and the next action that improves a validation objective. The action space can remain small: broad search, neighborhood exploration, relation relaxation, type tightening, rerank, answer, or restart. The reward can combine rank movement, recall gain, and budget cost. This formulation also makes personalization precise. A user state is not appended everywhere; it becomes one feature available to the policy, and the policy is rewarded only when using that feature improves retrieval.

\section{Conclusion}
\tb turns the common intuition that ``more explicit thinking before retrieval should help an adaptive KG retriever" into a measurable experiment. Under our controlled implementation, it does not. The unconditioned agent remains stronger on validation, held-out subsets, and paired shared-ID comparisons. The value of the study is that the reason becomes visible in the fact that repeated calls, zero-result calls, and misallocated exploration budget predict rank loss. This shifts the design target for personalized reasoned retrieval as the central problem is no longer how to add more profile text to the prompt, but how to learn a calibrated retrieval policy that knows when profile and trace information should influence the next action.
